% =========================================================================
% ACTA DE REUNIÓN — CS3081-006-2026 (BORRADOR)
% Plantilla de estilo basada en el modelo oficial de Google Docs del curso
% Compilar con:  tectonic Acta_Reunion6_BORRADOR.tex
% Regla del borrador: §3 (asistencia), §4 (tratado) y §6 (estados) van EN BLANCO.
% =========================================================================
\documentclass[11pt,a4paper]{article}

\usepackage[a4paper, top=2.6cm, bottom=2.4cm, left=2.5cm, right=2.5cm]{geometry}
\usepackage{fontspec}
\setmainfont{Carlito}                    % métrica idéntica a Calibri
\usepackage[table]{xcolor}
\usepackage{array}
\usepackage{longtable}
\usepackage{fancyhdr}
\usepackage{lastpage}
\usepackage{needspace}

% ------------------------- Paleta de la plantilla -------------------------
\definecolor{rojo}{HTML}{E4002B}   % marca CS3081
\definecolor{grisB}{HTML}{B7B7B7}  % bordes de tabla
\definecolor{grisH}{HTML}{202124}  % encabezados de tabla (negro)
\definecolor{grisL}{HTML}{F3F4F6}  % celdas de etiqueta
\definecolor{rosa}{HTML}{FCE8EC}   % caja PROPÓSITO / REGISTRO
\definecolor{grisV}{HTML}{555555}  % valores (itálica)
\definecolor{grisN}{HTML}{666666}  % notas y pie
\definecolor{tinto}{HTML}{333333}

\arrayrulecolor{grisB}
\renewcommand{\arraystretch}{1.55}
\setlength{\tabcolsep}{4pt}

% ------------------------- Encabezado y pie corridos ----------------------
\pagestyle{fancy}
\fancyhf{}
\fancyhead[L]{\footnotesize\textbf{\color{rojo}CS3081\hspace{0.7em}|\hspace{0.7em}INGENIERÍA DE SOFTWARE}\hspace{1.4em}{\color{grisN}·\hspace{0.7em}Acta de reunión con usuarios}}
\fancyfoot[R]{\footnotesize\color{grisN}Uso académico · Proyecto real\hspace{2em}Página \thepage\ de \pageref{LastPage}}
\renewcommand{\headrulewidth}{0pt}
\renewcommand{\footrulewidth}{0pt}
\setlength{\headheight}{14pt}

% ------------------------- Comandos de estilo -----------------------------
\newcommand{\asec}[1]{\Needspace*{13\baselineskip}\par\vspace{18pt}{\noindent\fontsize{14}{17}\selectfont\textbf{#1}}\par\vspace{9pt}}
\newcommand{\cajanota}[2]{\par\noindent\fcolorbox{grisB}{rosa}{\parbox{\dimexpr\linewidth-2\fboxsep-2\fboxrule}{\footnotesize \textbf{\textcolor{rojo}{#1: }}{\color{tinto}#2}}}\par\vspace{7pt}}
\newcommand{\hc}[1]{\cellcolor{grisH}\textcolor{white}{\textbf{#1}}}
\newcommand{\lb}[1]{\cellcolor{grisL}\textbf{#1}}
\newcommand{\vl}[1]{\textcolor{grisV}{\textit{#1}}}
\newcolumntype{L}[1]{>{\raggedright\arraybackslash}p{#1}}
\newcolumntype{C}[1]{>{\centering\arraybackslash}p{#1}}

\setlength{\parindent}{0pt}
\setlength{\parskip}{4pt}

\begin{document}

% ============================ TÍTULO ======================================
{\centering {\fontsize{22}{26}\selectfont\textbf{ACTA DE REUNION}}\par}
\vspace{10pt}

% ============================ 1. DATOS ====================================
\asec{1. DATOS GENERALES}
{\footnotesize
\begin{tabular}{|L{2.35cm}|L{5.45cm}|L{2.1cm}|L{4.3cm}|}
\hline
\lb{Código de acta} & \vl{CS3081-006-2026} & \lb{Fecha} & \vl{15/09/2026} \\ \hline
\lb{Hora de inicio} & \vl{10:00 a.m.} & \lb{Hora de término} & \vl{10:25 a.m.} \\ \hline
\lb{Proyecto} & \multicolumn{3}{L{12.0cm}|}{\vl{Proyecto Igualab --- Plataforma de analítica de sostenibilidad (RAG); línea base: Plan de Proyecto v1.2}} \\ \hline
\lb{Módulo/Fase} & \vl{Conformidad del PO sobre RN/RF/RNF (validadas por correo); presentación del \textbf{rediseño del mockup} con todas las funcionalidades; \textbf{entrega del avance completo del Análisis y Diseño (A\&D)}, a culminar en fecha posterior (lunes/martes o, de lograrse, el \textbf{viernes 18/09/2026})} & \lb{Modalidad} & \vl{[ ] Presencial \ [X] Virtual} \\ \hline
\end{tabular}}

% ============================ 2. OBJETIVO =================================
\asec{2. OBJETIVO Y RESULTADO ESPERADO}
{\footnotesize
\begin{tabular}{|L{2.9cm}|L{11.3cm}|}
\hline
\lb{Objetivo de la reunión} & \vl{Registrar la \textbf{conformidad del Product Owner (Oscar Baldeón)} con las \textbf{Reglas de Negocio (RN), los Requerimientos Funcionales (RF) y los Requerimientos No Funcionales (RNF)} remitidos por correo; \textbf{presentar al cliente el rediseño del mockup} con el conjunto completo de funcionalidades de la fase 1; y \textbf{presentar el avance completo del Análisis y Diseño (A\&D)} (REQ-13/14), con su numeración y trazabilidad, los casos de uso y el modelo de datos, \textbf{culminándolo en fecha posterior} (lunes/martes o, de lograrse, el viernes 18/09/2026).} \\ \hline
\lb{Resultado esperado} & \vl{\textbf{Conformidad del PO registrada} sobre RN/RF/RNF (evidencia: correo de conformidad); \textbf{rediseño del mockup presentado y aprobado} por el cliente; y \textbf{avance completo del Análisis y Diseño presentado/entregado} (RN 001--039, RF 001--056, RNF 001--036, casos de uso y modelo de datos), con su \textbf{culminación prevista en fecha posterior} (lunes/martes o, de lograrse, el viernes 18/09/2026).} \\ \hline
\end{tabular}}

% ============================ 3. PARTICIPANTES ============================
\newpage
\asec{3. PARTICIPANTES}
{\footnotesize\itshape\color{grisN} La asistencia se marca DURANTE la reunión.\par}
\vspace{4pt}
{\footnotesize
\begin{tabular}{|C{0.8cm}|L{4.1cm}|L{3.1cm}|L{4.4cm}|C{1.8cm}|}
\hline
\hc{N°} & \hc{Nombre y apellido} & \hc{Organización / equipo} & \hc{Rol en la reunión} & \hc{Asistencia} \\ \hline
1 & \vl{Oscar Rafael Baldeón Montoro} & \vl{Cliente} & \vl{Cliente / Product Owner} & \vl{[X] Sí \ [ ] No} \\ \hline
2 & \vl{\textbf{Fabricio Godofredo Ladera La Torre}} & \vl{Equipo UTEC} & \vl{\textbf{Jefe de Proyecto (PM)}} & \vl{[X] Sí \ [ ] No} \\ \hline
3 & \vl{Juan Renato Flores Pascual} & \vl{Equipo UTEC} & \vl{Desarrollador Frontend} & \vl{[ ] Sí \ [X] No} \\ \hline
4 & \vl{Carlos David Ordinola Ortega} & \vl{Equipo UTEC} & \vl{Analista Funcional} & \vl{[X] Sí \ [ ] No} \\ \hline
5 & \vl{Luis Javier Millones Carrasco} & \vl{Equipo UTEC} & \vl{Analista Funcional} & \vl{[X] Sí \ [ ] No} \\ \hline
6 & \vl{Juan Marcelo Ferreyra Gonzales} & \vl{Equipo UTEC} & \vl{Desarrollador Backend} & \vl{[ ] Sí \ [X] No} \\ \hline
7 & \vl{Mauricio Gabriel Gonzalez Kremer} & \vl{Equipo UTEC} & \vl{Tester} & \vl{[ ] Sí \ [X] No} \\ \hline
8 & \vl{Alonso Aarón Benites Camacho} & \vl{Equipo UTEC} & \vl{Tester} & \vl{[X] Sí \ [ ] No} \\ \hline
\end{tabular}}

% ============================ 4. AGENDA ===================================
\asec{4. AGENDA}
{\footnotesize\itshape\color{grisN} La columna ``Tratado'' se marca DURANTE la reunión.\par}
\vspace{4pt}
{\footnotesize
\begin{tabular}{|C{0.8cm}|L{8.6cm}|L{2.6cm}|C{2.2cm}|}
\hline
\hc{N°} & \hc{Tema} & \hc{Responsable} & \hc{Tratado} \\ \hline
1 & \vl{\textbf{Conformidad del PO (Oscar) sobre RN, RF y RNF}: validación de lo remitido y aprobado por correo} & \vl{Oscar / Fabricio (PM)} & \vl{[X] Sí \ [ ] No} \\ \hline
2 & \vl{\textbf{Presentación del rediseño del mockup} con el conjunto completo de funcionalidades de la fase 1} & \vl{Juan Renato (Frontend) / Equipo} & \vl{[X] Sí \ [ ] No} \\ \hline
3 & \vl{\textbf{Avance del Análisis y Diseño (A\&D)}: numeración y trazabilidad (RN/RF/RNF), casos de uso y modelo de datos (avance completo presentado/entregado; la \textbf{culminación se dará en fecha posterior}, lunes/martes o, de lograrse, el viernes 18/09)} & \vl{Luis / Carlos (Analistas)} & \vl{[X] Sí \ [ ] No} \\ \hline
\end{tabular}}

% ============================ 5. DESARROLLO ===============================
\newpage
\asec{5. DESARROLLO DE LA REUNIÓN}
{\footnotesize\itshape\color{grisN} Síntesis propuesta; se ajusta durante/tras la reunión.\par}
\vspace{4pt}
{\footnotesize
\begin{longtable}{|C{0.8cm}|L{3.0cm}|L{7.6cm}|L{2.8cm}|}
\hline
\rowcolor{grisH}\hc{N°} & \hc{Tema / necesidad} & \hc{Síntesis de lo conversado} & \hc{Conclusión / estado} \\ \hline
\endfirsthead
\hline
\rowcolor{grisH}\hc{N°} & \hc{Tema / necesidad} & \hc{Síntesis de lo conversado} & \hc{Conclusión / estado} \\ \hline
\endhead
1 & \vl{Conformidad del PO sobre RN/RF/RNF} & \vl{El equipo remitió al PO (Oscar) por correo el conjunto de \textbf{Reglas de Negocio (RN 001--039), Requerimientos Funcionales (RF 001--056) y Requerimientos No Funcionales (RNF 001--036)} para su revisión. \textbf{El PO confirmó que todo está conforme y sin observaciones} con lo indicado, quedando la conformidad registrada como evidencia (correo).} & \vl{\textbf{Conforme --- sin observaciones}} \\ \hline
2 & \vl{Presentación del rediseño del mockup} & \vl{Se muestra al cliente el \textbf{rediseño del mockup} de la plataforma, con el \textbf{conjunto completo de funcionalidades} de la fase 1 (login/RBAC de 2 roles, gestión de usuarios, ingesta de documentos \texttt{.md}, asistente IA/RAG con citado, revisión de brechas GRI y sanciones, generación y descarga de reportes PDF, auditoría). El rediseño se presenta alineado a las decisiones de alcance vigentes (sin módulo de Bolsa, sin LinkedIn y sin dashboards en fase 1).} & \vl{\textbf{Presentado y aprobado por el PO}} \\ \hline
3 & \vl{Avance del Análisis y Diseño (A\&D)} & \vl{El equipo presenta y \textbf{entrega el avance completo} del \textbf{Análisis y Diseño}, con la numeración y trazabilidad de \textbf{RN, RF y RNF}, los \textbf{casos de uso} y el \textbf{modelo de datos}, validando con el cliente que el documento refleja lo acordado en las actas 4 y 5. \textbf{La culminación del A\&D (documento acabado) queda prevista en fecha posterior} (lunes/martes o, de lograrse, la reunión del viernes 18/09/2026), sin comprometer el cierre a esa fecha.} & \vl{\textbf{Avance completo presentado --- culminación posterior}} \\ \hline
\end{longtable}}

% ============================ 6. REQUISITOS ===============================
\asec{6. NECESIDADES, REQUISITOS Y CAMBIOS IDENTIFICADOS}
{\footnotesize\itshape\color{grisN} Registre solo lo comprendido en la reunión. La columna ``Estado'' se completa DURANTE/TRAS la reunión.\par}
\vspace{5pt}
{\footnotesize
\begin{longtable}{|C{1.3cm}|L{5.3cm}|C{1.3cm}|C{1.3cm}|L{3.4cm}|C{1.6cm}|}
\hline
\rowcolor{grisH}\hc{ID} & \hc{Necesidad / requisito / cambio} & \hc{Tipo} & \hc{Prioridad} & \hc{Criterio o evidencia} & \hc{Estado} \\ \hline
\endfirsthead
\hline
\rowcolor{grisH}\hc{ID} & \hc{Necesidad / requisito / cambio} & \hc{Tipo} & \hc{Prioridad} & \hc{Criterio o evidencia} & \hc{Estado} \\ \hline
\endhead
REQ-21 & \vl{\textbf{Conformidad del PO (Oscar) con las RN, RF y RNF} remitidos por correo: el cliente revisó el conjunto normativo de la fase 1 (\textbf{RN 001--039, RF 001--056, RNF 001--036}) y manifestó su conformidad \textbf{sin observaciones}. Se registra como validación formal de los requerimientos de la fase 1.} & \vl{Documento} & \vl{Alta} & \vl{Correo de conformidad del PO (Oscar) sobre RN/RF/RNF} & \vl{\textbf{Aprobado}} \\ \hline
REQ-22 & \vl{\textbf{Presentación y aprobación del rediseño del mockup}: se presentó al cliente el rediseño del prototipo con el \textbf{conjunto completo de funcionalidades} de la fase 1, alineado a las decisiones de alcance (2 roles, sin Bolsa, sin LinkedIn, sin dashboards en fase 1).} & \vl{Documento} & \vl{Alta} & \vl{Rediseño presentado al PO; conformidad registrada en esta acta} & \vl{\textbf{Aprobado}} \\ \hline
REQ-23 & \vl{\textbf{Entrega/presentación del avance completo del Análisis y Diseño (A\&D)}: se registra la \textbf{entrega del avance completo} del A\&D en esta reunión (numeración y trazabilidad de RN/RF/RNF, casos de uso y modelo de datos) y se acuerda \textbf{culminar el documento (acabado) en fecha posterior} (lunes/martes o, de lograrse, el viernes 18/09/2026), sin comprometer el cierre a esa fecha.} & \vl{Documento} & \vl{Alta} & \vl{Avance completo presentado/entregado; culminación en fecha posterior (lunes/martes o viernes 18/09)} & \vl{\textbf{Aprobado}} \\ \hline
\end{longtable}}

% ============================ 7. ACUERDOS =================================
\asec{7. ACUERDOS Y COMPROMISOS}
{\footnotesize\itshape\color{grisN} Cada compromiso debe expresar una acción verificable, un responsable único y una fecha acordada. Se completa durante/tras la reunión.\par}
\vspace{5pt}
{\footnotesize
\begin{longtable}{|C{0.8cm}|L{6.3cm}|L{2.5cm}|C{1.8cm}|C{1.5cm}|L{1.3cm}|}
\hline
\rowcolor{grisH}\hc{N°} & \hc{Acción / entregable} & \hc{Responsable} & \hc{Fecha límite} & \hc{Estado} & \hc{Evidencia} \\ \hline
\endfirsthead
\hline
\rowcolor{grisH}\hc{N°} & \hc{Acción / entregable} & \hc{Responsable} & \hc{Fecha límite} & \hc{Estado} & \hc{Evidencia} \\ \hline
\endhead
1 & \vl{\textbf{Conformidad del PO sobre RN/RF/RNF registrada}: el cliente manifiesta su conformidad sin observaciones con el conjunto normativo remitido por correo} & \vl{Oscar (Cliente)} & \vl{15/09/2026 (en reunión)} & \vl{\textbf{Aprobado}} & \vl{Esta acta / correo} \\ \hline
2 & \vl{\textbf{Rediseño del mockup presentado y aprobado} por el cliente, con el conjunto completo de funcionalidades de la fase 1} & \vl{Juan Renato (Frontend)} & \vl{15/09/2026 (en reunión)} & \vl{\textbf{Aprobado}} & \vl{Esta acta / demo} \\ \hline
3 & \vl{\textbf{Entregar/presentar el avance completo del Análisis y Diseño (A\&D)}, con su numeración y trazabilidad, los casos de uso y el modelo de datos, y \textbf{culminar el documento en fecha posterior} (lunes/martes o, de lograrse, el viernes)} & \vl{Luis / Carlos (Analistas)} & \vl{18/09/2026 (avance completo)} & \vl{\textbf{Aprobado}} & \vl{A\&D V1.2} \\ \hline
4 & \vl{\textbf{Enviar el acta 6 por correo}} & \vl{Fabricio (PM)} & \vl{16/09/2026} & \vl{\textbf{Realizado}} & \vl{Correo} \\ \hline
\end{longtable}}

% ============================ 8. PRÓXIMA ==================================
\asec{8. PRÓXIMA INTERACCIÓN Y SEGUIMIENTO}
{\footnotesize
\begin{tabular}{|L{2.5cm}|L{4.6cm}|L{2.6cm}|L{4.5cm}|}
\hline
\lb{Próxima reunión} & \vl{\textbf{18/09/2026 (viernes)} --- cadencia de \textbf{2 reuniones por semana}} & \lb{Objetivo} & \vl{Presentar el \textbf{avance completo del Análisis y Diseño} y validar su \textbf{culminación} (lunes/martes o, de lograrse, en esa misma reunión), con el alcance de fase 1 conforme a lo aprobado} \\ \hline
\lb{Canal de seguimiento} & \vl{WhatsApp (coordinación diaria) / Correo (actas)} & \lb{Responsable del acta} & \vl{Fabricio Godofredo Ladera La Torre (PM)} \\ \hline
\lb{Fecha de envío del acta} & \vl{16/09/2026} & \lb{Ubicación del archivo} & \vl{Repositorio del equipo CS3081} \\ \hline
\end{tabular}}

% ============================ 9. VALIDACIÓN ===============================
\asec{9. VALIDACIÓN DEL ACTA/FIRMAS}
{\footnotesize\itshape\color{grisN} La validación puede realizarse mediante firma, correo, mensaje en el canal acordado o confirmación registrada. Adjunte la evidencia correspondiente.\par}
\vspace{5pt}
{\footnotesize
\begin{tabular}{|L{4.7cm}|L{4.7cm}|L{4.7cm}|}
\hline
\hc{Representante del usuario / cliente} & \hc{Representante del equipo} & \hc{Docente o ACL (si corresponde)} \\ \hline
\vl{Nombre: Oscar Rafael Baldeón Montoro\newline Cargo / rol: Cliente / Product Owner} & \vl{Nombre: Fabricio Godofredo Ladera La Torre\newline Cargo / rol: Jefe de Proyecto (PM)} & \vl{Nombre: Teofilo Chambilla Aquino\newline Cargo / rol: Docente (acompañamiento)} \\ \hline
\vl{Método de validación: Firma y correo} & \vl{Método de validación: Correo (envío del acta)} & \vl{Método de validación: [Firma / correo]} \\ \hline
\vl{Fecha: 15/09/2026\newline Observaciones: Conforme con el plan v1.2, las decisiones de alcance y el formato de ingesta .md} & \vl{Fecha: 15/09/2026\newline Observaciones: Acta enviada por correo} & \vl{Fecha: [dd/mm/aaaa]\newline Observaciones: [Completar]} \\ \hline
\end{tabular}}

% ============================ 10. APROBACIÓN ==============================
\asec{10. APROBACIÓN (V°B° CEO)}
El presente acta es revisada y aprobada por la Gerencia General antes de su envío al cliente.
\vspace{4pt}

{\itshape Oscar Rafael Baldeón Montoro}
\vspace{10pt}

\textbf{Nombre:} Oscar Rafael Baldeón Montoro

\textbf{Cargo:} CEO / Gerente General

\vspace{28pt}

\textbf{Firma:} \rule[-2pt]{6.5cm}{0.4pt}

\textbf{Fecha:} [dd/mm/aaaa]

\vspace{6pt}
\label{LastPage}
\end{document}
